% ===========================================================================
\chapter{Introdução e Fundamentação do Design Centrado no Usuário}
\label{cap_introducao_fundamentacao}
% ===========================================================================

A formação acadêmica de profissionais de saúde no Brasil enfrenta um desafio pedagógico de proporções epidemiológicas: o ensino sistemático da identificação precoce, análise causal e mitigação de incidentes assistenciais e eventos adversos. Historicamente, o ensino hospitalar operou sob paradigmas punitivos, nos quais falhas assistenciais eram atribuídas exclusivamente à culpa individual de enfermeiros ou médicos, inibindo notificações espontâneas e perpetuando ciclos de negligência sistêmica.

A metodologia \textit{Global Trigger Tool} (GTT), concebida pelo \textit{Institute for Healthcare Improvement} \cite{ihi2009}, revolucionou a vigilância clínica internacional ao instituir um método retrospectivo baseado no rastreamento de ``gatilhos'' (pistas textuais ou laboratoriais em prontuários que indicam a ocorrência de dano). O método IHI preconiza que a auditoria clínica de cada prontuário seja concluída em no máximo vinte minutos, exigindo que o auditor desenvolva foco investigativo de alta densidade sem perder-se em minúcias irrelevantes.

Contudo, a transposição dessa metodologia para o ambiente educacional universitário -- especificamente na Universidade Federal de Sergipe (UFS), em cooperação multidisciplinar entre o Departamento de Computação (DCOMP) e o Departamento de Enfermagem -- esbarra em barreiras cognitivas severas. A execução de auditorias práticas por meio de fichas físicas em papel ou planilhas eletrônicas desarticuladas resulta em extrema sobrecarga cognitiva, ansiedade temporal, fadiga visual e ausência de consolidação analítica de indicadores docentes.

Nesse cenário, o desenvolvimento do SIGEA-GTT (Sistema Inteligente de Gestão de Eventos Adversos -- Global Trigger Tool) apoia-se nos preceitos do Design Centrado no Usuário (DCU). Este documento formaliza o processo de engenharia de usabilidade, detalhando as personas arquetípicas, os mapas de empatia em seis quadrantes e as jornadas dos usuários que alicerçam as decisões de interface e arquitetura de software da plataforma.

% ===========================================================================
\section{A Norma ISO 9241-210 e o Paradigma do DCU}
\label{sec_norma_iso_9241}
% ===========================================================================

A norma internacional ISO 9241-210:2019 \cite{iso9241} estabelece os princípios orientadores para o projeto centrado no ser humano em sistemas interativos de computação. De acordo com a norma, o projeto centrado no usuário não é uma etapa periférica ou decorativa de desenho de telas, mas uma abordagem de engenharia que visa otimizar a eficácia, a eficiência e a satisfação do operador humano, prevenindo danos à saúde e minimizando o erro humano em contextos críticos.

A ISO 9241-210 estrutura o desenvolvimento iterativo em quatro atividades essenciais:
\begin{enumerate}
    \item \textbf{Compreensão e Especificação do Contexto de Uso}: Identificação das características dos operadores, tarefas clínicas executadas, ferramentas pedagógicas e ambientes organizacionais;
    \item \textbf{Especificação dos Requisitos do Usuário}: Tradução das necessidades humanas, dores cognitivas e restrições ergonômicas em metas explícitas de interação;
    \item \textbf{Concepção de Soluções de Design}: Criação de arquitetura de informação, fluxos de navegação, modelos mentais, protótipos de baixa e alta fidelidade;
    \item \textbf{Avaliação Baseada no Usuário}: Validação de usabilidade com usuários representativos frente a critérios quantitativos e qualitativos preestabelecidos.
\end{enumerate}

No SIGEA-GTT, a aplicação desse ciclo iterativo assegura que os três perfis autorizados de acesso (Discentes, Docentes e Administradores) encontrem uma ferramenta ergonomicamente adaptada às suas rotinas cognitivas no ambiente dos laboratórios da UFS.

% ===========================================================================
\section{Metodologia Goal-Directed Design e Teoria de Personas}
\label{sec_teoria_personas}
% ===========================================================================

A metodologia \textit{Goal-Directed Design} (GDD), formulada por Alan Cooper \cite{cooper1999, cooper2014}, postula que sistemas digitais frequentemente falham não por ausência de capacidade computacional, mas por tentarem satisfazer um ``usuário elástico'' -- uma abstração fictícia cujas habilidades e desejos são continuamente distorcidos pela equipe técnica durante a implementação.

Para superar esse vício de projeto, Cooper introduziu o conceito de \textbf{Personas}: arquétipos realistas e compostos, derivados de observações empíricas e entrevistas com usuários reais, que personificam objetivos, atitudes, modelos mentais, limitações e contextos de trabalho bem delimitados \cite{cooper2014}. 

As personas distinguem-se em diferentes categorias hierárquicas de projeto:
\begin{itemize}
    \item \textbf{Persona Primária}: O foco principal do design de determinada interface. Suas necessidades não podem ser atendidas por interfaces projetadas para outras personas sem gerar frustração inaceitável;
    \item \textbf{Persona Secundária}: Perfil cujas demandas são majoritariamente contempladas pelo design da persona primária, demandando apenas ajustes ou extensões pontuais;
    \item \textbf{Persona Terciária}: Usuário de suporte ou governança técnica cujas operações administrativas operam em segundo plano;
    \item \textbf{Anti-Persona}: Perfil explicitamente excluído do escopo do sistema, delimitando o que a solução \textbf{não} deve ser e evitando dispersão de esforço de engenharia.
\end{itemize}

% ===========================================================================
\section{O Empathy Map Canvas de Gray e Osterwalder}
\label{sec_empathy_map_canvas}
% ===========================================================================

O Mapa de Empatia (\textit{Empathy Map Canvas}), desenvolvido por Dave Gray \cite{gray2010} no âmbito da XPLANE e consolidado por Osterwalder e Pigneur \cite{osterwalder2014}, é um artefato visual de síntese cognitiva que transporta a equipe de engenharia para o interior do universo perceptual do usuário.

A ferramenta estrutura o modelo sensorial e emocional da pessoa em seis quadrantes interdependentes:
\begin{enumerate}
    \item \textbf{O que pensa e sente?}: As aspirações mais profundas, preocupações reais, sentimentos de ansiedade, inseguranças profissionais e motivações intrínsecas;
    \item \textbf{O que escuta?}: As mensagens transmitidas por pares, professores, supervisores hospitalares, diretrizes de conselhos profissionais e a cultura institucional;
    \item \textbf{O que vê?}: Os estímulos visuais cotidianos, as interfaces de prontuários com as quais interage, a sobrecarga de papéis e o ambiente físico do laboratório ou enfermaria;
    \item \textbf{O que fala e faz?}: A linguagem e comportamento exteriorizados em público, termos técnicos empregados, posturas em sala de aula e atitudes de busca;
    \item \textbf{Dores (\textit{Pains})}: As frustrações, medos de cometer erros, perdas de tempo, atritos com softwares obsoletos e cansaço visual;
    \item \textbf{Ganhos (\textit{Gains})}: Os critérios de sucesso, sentimentos de competência adquirida, economia de tempo, clareza pedagógica e reconhecimento acadêmico.
\end{enumerate}

% ===========================================================================
\section{Design Emocional de Norman e Engenharia de Usabilidade de Nielsen}
\label{sec_norman_nielsen}
% ===========================================================================

O design de software voltado para cenários clínicos críticos deve harmonizar o rigor funcional com a ressonância psicológica. Donald Norman \cite{norman2004, norman2013} conceitua três níveis cognitivos de processamento no cérebro humano:
\begin{itemize}
    \item \textbf{Nível Visceral}: Resposta emocional imediata e pré-consciente aos aspectos estéticos, paleta de cores, harmonia tipográfica e clareza visual;
    \item \textbf{Nível Comportamental}: Experiência prática de uso cotidiano, focada em eficácia, tempo de resposta da interface, inteligibilidade de \textit{affordances} e ausência de atrito;
    \item \textbf{Nível Reflexivo}: Memória consolidada, atribuição de significado e orgulho pessoal associados ao domínio de uma ferramenta que torna o usuário um profissional mais qualificado.
\end{itemize}

Complementarmente, Jakob Nielsen \cite{nielsen1993, nielsen1994} formalizou os critérios objetivos de Engenharia de Usabilidade, definindo atributos quantificáveis de facilidade de aprendizado (\textit{learnability}), eficiência de uso, memorabilidade, baixas taxas de erro e satisfação subjetiva, além das consagradas dez heurísticas de inspeção que orientam o projeto e a verificação das telas do SIGEA-GTT.

A Tabela \ref{tab_comparativo_papeis_ucd} sumariza a correspondência entre os papéis de usuários, objetivos pedagógicos e desafios cognitivos mapeados no projeto.

\begin{table}[!ht]
\centering
\caption{Matriz Comparativa dos Atores do SIGEA-GTT sob a Ótica do DCU}
\label{tab_comparativo_papeis_ucd}
\footnotesize
\begin{tabularx}{\textwidth}{p{2.6cm} p{2.4cm} p{4.2cm} X}
\toprule
\textbf{Ator do Sistema} & \textbf{Papel no DCU} & \textbf{Objetivo Central de Interação} & \textbf{Principal Desafio Cognitivo e Ergonômico} \\
\midrule
\textbf{Discente} & Persona Primária & Aprender auditoria retrospectiva GTT, rastrear gatilhos e aplicar ferramentas da qualidade. & Gerenciar a pressão temporal dos 20 minutos de auditoria e superar a ansiedade clínica. \\
\addlinespace
\textbf{Docente} & Persona Secundária & Gerenciar turmas, elaborar prontuários simulados e avaliar com gabarito comparativo. & Eliminar a fadiga na correção manual de dezenas de papéis e consolidar taxas do IHI. \\
\addlinespace
\textbf{Administrador} & Persona Terciária & Garantir integridade de acessos institucionais e parametrizar catálogo de gatilhos. & Manter conformidade com a LGPD e governança estrita de perfis sem gerar atrito técnico. \\
\addlinespace
\textbf{Auditor Punitivo} & Anti-Persona & Vigiar ou penalizar profissionais de saúde em prontuários assistenciais reais. & Incompatível com o propósito pedagógico formativo e não-punitivo do SIGEA-GTT. \\
\bottomrule
\end{tabularx}
\fonte{Elaborado pelo autor (2026).}
\end{table}
